\chapter{Change-of-base, $\Fam(\mathcal{C}^{\mathrm{op}})$ generality, and the $\{\mathrm{Id}\}$ classification}
\label{ch:cob}
How far does the equation generalise? The natural setting is containers over a
general base, i.e.\ $\Fam(\mathcal{C}^{\mathrm{op}})$, with change-of-base along
functors and monads. This chapter assembles the positive generality results
($m$-containers, codensity/polynomiality, the four monoidal structures) and the
sharp negative: the self-enriched full-and-faithful change-of-base container
semantics collapses to $\{\mathrm{Id}\}$ --- no nontrivial effect monad supports
it. Knowing the boundary is itself the result.

%========================================================================
\section{The base as a degree of freedom}
\label{sec:cob-setup}

Recall from Chapter~\ref{ch:prelim} that a container is an object of
$\Fam(\Set^{\op})$ (Notation~\ref{not:prelim-cont}) and that its extension is a
left Kan extension along the representable embedding (Theorem~\ref{thm:prelim-lan}),
\[
  \den{S,P}\,X \;=\; \sum_{s\in S}\Set(P_s,X)\;=\;\sum_{s\in S}X^{P_s}.
\]
The formula has exactly one silent parameter: the hom-set $\Set(P_s,X)$. Replace
$\Set$ by a category $\mathcal{C}$ and read positions out in $\mathcal{C}$'s
internal hom, and the same syntax presents a wholly different world.

\begin{definition}[Containers over a base]\label{def:cob-fam}
For a closed symmetric monoidal category $(\mathcal{C},\otimes,I,[-,-])$ with
small coproducts, $\Fam(\mathcal{C}^{\op})$ has objects $(S,P)$ with $S$ a set of
shapes and $P_s\in\mathcal{C}$ a \emph{position object} for each $s$, and
morphisms $(S,P)\to(T,Q)$ a shape map $f\colon S\to T$ together with, per $s$, a
$\mathcal{C}$-map $Q_{f(s)}\to P_s$ (contravariant in positions). The
\emph{extension} into $\mathcal{C}$-enriched endofunctors is
\[
  \den{S,P}\;=\;\sum_{s\in S}[P_s,-]\;\colon\;\mathcal{C}\to\mathcal{C}.
\]
\end{definition}

A single construction thus covers the whole programme, selected by the base:
$\mathcal{C}=\Set$ gives ordinary containers; $\mathcal{C}=\mathrm{Kl}(M)$ for a
monad $M$ gives \emph{effectful} ($M$-)containers, whose positions are read in a
Kleisli category; $\mathcal{C}=\VecC$ gives linear containers; and
$\mathcal{C}=\Cat$ recovers directed containers and hence small categories
(Chapter~\ref{ch:prelim}, \S\ref{sec:prelim-dcont}). The assignment
$\mathcal{C}\mapsto\Fam(\mathcal{C}^{\op})$ is functorial: a functor
$F\colon\mathcal{C}\to\mathcal{D}$ induces $\Fam(F^{\op})$, acting as $F$ on
positions and identically on shapes, and the total category assembling all bases
is a bifibration over $\Set$ whose fibrewise opposite is the contravariance in
positions \cite{vonGlehn}. \emph{Change of base} along $F\colon\Set\to\mathcal{C}$
is the transport of the standard theory into $\mathcal{C}$; when $F$ is
(co)monoidal it carries (co)monoidal structure, and that is where effects and
directedness enter.

The remainder of the chapter is a suite of positive generality results
(\S\S\ref{sec:cob-mcont}--\ref{sec:cob-admiss}), a diagnosis of the linear wall
(\S\ref{sec:cob-vec}), and the collapse that fixes the exact boundary of a
\emph{uniform} enriched semantics (\S\ref{sec:cob-collapse}).

%========================================================================
\section{Positive suite I: $m$-containers, codensity, polynomiality}
\label{sec:cob-mcont}

The Kleisli base is the central instance and the one where two Kan-theoretic
invariants of the effect first appear. Fix a monad $M=(M,\eta,\mu)$ on $\Set$; an
\emph{$M$-container} is a container $(S,P)$ whose positions are read in
$\mathrm{Kl}(M)$, with $M$-extension $\den{S,P}_M X=\sum_{s}\Set(P_s,MX)$. All
results of this section are \emph{peer-reviewed} (Rick).

\begin{theorem}[Identification and factorisation; \emph{peer-reviewed}]\label{thm:cob-mfactor}
$M\text{-}\Cont=\Fam(\mathrm{Kl}(M)^{\op})$, and as endofunctors, naturally in
$(S,P)$,
\[
  \den{S,P}_M\;=\;\den{S,P}\circ M .
\]
\end{theorem}

So an $M$-container functor is an \emph{ordinary polynomial functor with a
trailing $M$}. This one factorisation is the lever for everything below: every
structural question about $\den{-}_M$ becomes a question about the single
trailing $M$. The proof is a direct unfolding of the two hom-descriptions.

\begin{theorem}[Faithful always; fullness $=$ codensity; \emph{peer-reviewed}]\label{thm:cob-codensity}
$\den{-}_M\colon M\text{-}\Cont\to[\Set,\Set]$ is \emph{always faithful}
(the Kleisli unit $\eta$ is a section, so container data is read back off any
induced transformation). It is \emph{full} --- hence fully faithful --- if and
only if \textup{(i)} each functor $\Set(A,M{-})$ is connected, and \textup{(ii)}
$M$ is \emph{codense}, i.e.\ the canonical monad map $M\Rightarrow\Ran_M M$ into
its own codensity monad is an isomorphism.
If $M$ is affine ($M1\cong1$) then \textup{(i)} holds automatically, so
\emph{affine $+$ codense $\Rightarrow$ fully faithful}.
\end{theorem}

The engine is that precomposition $(-)\circ M$ has right adjoint $\Ran_M$, so
every hom-set $\Nat(\Set(A,M{-}),H)$ is a value of the right Kan extension along
$M$; the fullness obstruction is exactly the failure of the unit
$M\Rightarrow\Ran_M M=\Phi$ (the codensity monad, $\Phi(A)=\Nat(\Set(A,M{-}),M)$)
to be invertible. The verdict is genuinely selective.

\begin{example}[The effects sort themselves]\label{ex:cob-effects}
The finite-distribution monad $\mathsf{D}$ is affine and codense, so
$\den{-}_{\mathsf{D}}$ \emph{is} fully faithful (modulo the standard rigidity
$\Nat(\mathsf{D},\mathsf{D})=\{\id\}$). By contrast $\mathsf{Maybe}$, exceptions $X+E$, the writer
$E\times(-)$ and the reader $(-)^R$ ($|R|\ge2$) all fail codensity and so lose
fullness. The reader is affine but \emph{not} codense
($\Phi(A)=(RA)^R\neq A^R=MA$), showing (i) and (ii) are independent conditions.
\end{example}

\begin{remark}[The wrong criterion]\label{rem:cob-wrongcrit}
The natural guess ``$\den{-}_M$ is full iff $M$ preserves coproducts'' is false
for the target $[\Set,\Set]$: that is the criterion for the coarser
Kleisli-enriched target, and it gives the wrong verdict on $\mathsf{D}$. The writer
$E\times(-)$ \emph{preserves} coproducts yet fails codensity
($\Phi(1)=|E|^{|E|}\neq|E|$ for $|E|\ge2$), a clean refutation. The controlling
invariant is the codensity monad $\Ran_M M$, not coproduct preservation.
\end{remark}

The second invariant governs composition. Call $M$ \emph{polynomial} if
$M\cong\sum_{i\in I}(-)^{B_i}$, i.e.\ $M=\den{\lceil M\rceil}$ for a container
$\lceil M\rceil$.

\begin{theorem}[Composition $=$ polynomiality; \emph{peer-reviewed}]\label{thm:cob-compose}
\emph{(Sufficiency, unconditional.)} If $M$ is polynomial then
$p\tri_M q:=p\tri\lceil M\rceil\tri q$ satisfies
$\den{p\tri_M q}_M\cong\den{p}_M\circ\den{q}_M$ naturally, so the essential image
of $\den{-}_M$ is closed under composition. The operation $\tri_M$ is associative
but has a two-sided unit up to container isomorphism \emph{only} when $M=\Id$: it
is semigroupal, not monoidal.
\emph{(Necessity, conditional.)} The converse --- closure $\Rightarrow M$
polynomial --- holds on an exactly delimited class strictly larger than the
\emph{fully symmetric} world: for super-polynomial $M$ (a cardinality law, ruling
out powerset, $\mathsf{D}$, ultrafilter, continuation); for analytic $M$ with
$M\emptyset=\emptyset$ (right-cancellation of plethysm, at unbounded degree and
nesting); for analytic $M$ of bounded wreath depth (an imprimitivity invariant
subsuming multiset and symmetric powers); and for the free-algebra monad of every
unital non-flat operad with a commutative binary composite and
$\Delta_{C_2}\notin\mathcal{S}$. It \emph{fails} for partially symmetric unital
operads: a ternary operation with cyclic ($C_3$) slot-symmetry and an absorbed
unit is non-polynomial yet composition-closed.
\end{theorem}

The one genuinely open point is the converse of the last clause
($\Delta_{C_2}\in\mathcal{S}\Rightarrow$ closed), verified for the all-binary
inhabitant at small arities and open in general. Within the fully symmetric world
--- which contains all the applied effect monads (multiset, distribution,
powerset) --- composition $\iff$ flatness $\iff$ polynomiality.

\begin{corollary}[The trade-off; \emph{peer-reviewed}]\label{cor:cob-tradeoff}
Fullness (codensity) and composability (polynomiality) are independent invariants
of $M$ that trade oppositely across the standard effects. The writer $E\times(-)$
is polynomial, hence composes, but is not codense, hence not full. The
distribution monad $\mathsf{D}$ is affine and codense, hence full, but is not polynomial,
hence does not compose. \emph{Probability keeps faithful naming and loses
composition; error keeps composition and loses faithful naming.} In particular
the folklore ``composes $\iff$ commutative and affine'' is false: $\mathsf{D}$ is affine
and commutative but does not compose, while $\mathsf{Maybe}$ is neither yet does.
\end{corollary}

%========================================================================
\section{Positive suite II: the change-of-base census}
\label{sec:cob-changebase}

Beyond the Kleisli identification, the same frame $\Fam((-)^{\op})$ absorbs a
census of monad/comonad transports, each a change of base or a fibrewise transport
read in the appropriate variance. All are \emph{proved} in the companion material.

\begin{proposition}[Position transfer: monads become comonads; \emph{proved}]\label{prop:cob-transfer}
Because positions are contravariant, a monad $M$ on $\Set$ postcomposed on the
position family, $G(S,P)=(S,M\circ P)$, is a \emph{comonad} on $\Cont$ --- with
counit $\eta$ and comultiplication $\mu$ on the fibres --- and this is a
biconditional: the comonad laws for $G$ hold iff the monad laws for $M$ do
(single-shape containers realise every fibre). Globally $G=\Lan_{(S,P)}M$. Dually
a comonad on $\Set$ yields a monad on $\Cont$. This is the slogan of the whole
frame in miniature: \emph{a monad on the base is a comonad on the fibrewise
opposite}.
\end{proposition}

\begin{proposition}[The $\forall/\exists$ liftings; \emph{proved}]\label{prop:cob-liftings}
Quantifying a predicate container over positions gives two liftings. The
$\exists$-lifting is exactly composition, $E=\tri$, a bifunctor on all of $\Cont$
with unit $\yy$. The $\forall$-lifting $A=\Pi$ is a bifunctor only on
\emph{cartesian} morphisms in its first argument (the pushforward is indexed by
set-sections of the backward position map, which are unique iff the map is a
bijection), and on that domain it is a left module of $(\Cont,\tri,\yy)$. The
mixed comparison $A\,X\,(E\,Y\,C)\to E\,(A\,X\,Y)\,C$ is an isomorphism iff $X$ is
\emph{linear}. This is the semantic content of $M$ lifting to containers iff $M$
is cartesian, seen one level down.
\end{proposition}

\begin{proposition}[The free monad on a container; \emph{proved}]\label{prop:cob-free}
A signature is a container, and its free monad --- the type of effect trees --- is
a monad on $\Cont$ built by iterated $\tri$-composition. In the linear instance
over a monoidal base, the free $\tri$-monoid on the one-shape container
$C=(\star,A)$ is $C^{*}=(\N,(A^{\otimes n})_n)$, with
$\den{C^{*}}(W)=\bigoplus_{n\ge0}A^{\otimes n}\otimes W=T(A)\otimes W$ the tensor
algebra: over $\VecC$ this is O'Neill's free monad, realised functorially.
\end{proposition}

The composition question of \S\ref{sec:cob-mcont} is the fifth census entry, and
the position transfer and liftings are the fibrewise ones; naming the frame turns
a list of theorems into a single diagram.

%========================================================================
\section{Positive suite III: the four monoidal structures}
\label{sec:cob-four}

The base $\Cont$ itself carries four monoidal structures, previewed in
Chapter~\ref{ch:prelim}, \S\ref{sec:prelim-day}: the coproduct $+$, the product
$\times$, the Dirichlet tensor $\Day$ (Definition~\ref{def:prelim-day}; on
containers the shapes and positions both multiply, cf.\
Proposition~\ref{prop:prelim-degen}), and the substitution operator $\tri$. Their
full classification is a topic in its own right; we record here only the map and
the closure results the rest of the dossier uses. All are \emph{proved}; the Lean
certifications are noted where they exist.

\paragraph{Why exactly these four.} Every monoidal structure $(\star,I)$ on $\Set$
induces one on $\Cont$ by Day convolution \cite{day1970,NiuSpivak}: with
$p\,\Day_\star\,q=(S_p\times S_q,(s,t)\mapsto p_s\star q_t)$ and unit $y^{I}$.
Call a structure \emph{convolutional} if it preserves coproducts in each variable
and closes the representables.

\begin{theorem}[Day classification; \emph{proved}]\label{thm:cob-dayclass}
$(\star,I)\mapsto(\Day_\star,y^{I})$ is an \emph{equivalence} from monoidal
structures on $\Set$ to convolutional monoidal structures on $\Cont$. Two of the
four are its values, $\times=\Day_{(+,\varnothing)}$ and
$\Day=\Day_{(\times,1)}$; the coproduct $+$ and the substitution $\tri$ lie
outside the family (they fail coproduct-preservation in a variable). This
explains the ``upside-down'' arithmetic --- the product's positions are a
\emph{disjoint union} because it is the Day convolution of the \emph{coproduct} of
sets, $y^{A+B}\cong y^{A}\times y^{B}$.
\end{theorem}

\begin{theorem}[Pointwise $=$ product; \emph{proved}]\label{thm:cob-pointwise}
Among \emph{all} monoidal structures on $\Cont$ --- Day or not --- the categorical
product $\times$ is the unique one whose extension is the pointwise product of
functors, $\den{p\ast q}\cong\den{p}\times\den{q}$. Being pointwise is the rare
property, not the generic one; in particular $\den{-}$ is strong monoidal into
$([\Set,\Set],\times)$ \emph{only} for the cartesian structure.
\end{theorem}

\begin{theorem}[The comparitor; \emph{proved}]\label{thm:cob-comparitor}
For each $p$, $p\,\Day\,(-)$ is the left Kan extension of $p\tri(-)$ along the
representable embedding, and the comparison $p\,\Day\,q\to p\tri q$ is the counit
of the resulting coreflection: $p\,\Day\,(-)$ is the terminal
coproduct-preserving approximation to $p\tri(-)$. The Dirichlet tensor is the
substitution operator with its dependency switched off --- and not by fiat: it is
the Day-ification of $\tri$, which is why the two agree on representables
($y^{A}\,\Day\,y^{B}=y^{A\times B}=y^{A}\tri y^{B}$).
\end{theorem}

\paragraph{Closures.} The Dirichlet tensor is closed, with internal hom the
container whose \emph{shapes are the container morphisms} \cite{NiuSpivak}; this
closure is machine-checked --- $(\Cont,\Day,\yy)$ is a closed monoidal category in
Lean~4, currying and uncurrying inverse on the nose (\emph{lean-verified},
\texttt{Dirichlet.lean}). The substitution operator $\tri$ is \emph{not}
left-closed but carries a right coclosure, whose comonoids are the directed
containers, hence the small categories \cite{AU16}. And $\Cont$ is
cartesian closed but \emph{not} locally cartesian closed, so it does not model
dependent type theory --- a fact worth stating plainly, and one the frontier
chapter revisits.

\paragraph{Interaction.} The pair $(\Day,\tri)$ is a normal duoidal category
\cite{SpivakSrinivasan}, and its interchange law houses the comparitor. Its double
comonoids --- objects that are comonoids for both tensors compatibly --- are
classified: they are exactly the \emph{sets of commutative monoids}, by a
fibrewise Eckmann--Hilton collapse (\emph{proved}), correcting the plausible guess
``the degenerate polynomials'' in both directions.

%========================================================================
\section{Positive suite IV: when does substitution exist over a base?}
\label{sec:cob-admiss}

The Dirichlet tensor $p\Day q=(S\times T,(P_s\otimes Q_t))$ exists over
\emph{every} closed monoidal base. The substitution product $\tri$ --- the one
whose comonoids are categories --- is different: it mostly does not exist, and
where it does, its existence is not a matter of degree but a sharp theorem. Call
$\mathcal{C}$ \emph{$\tri$-admissible} if the image of $\den{-}$ is closed under
composition, $\den{p\tri q}\cong\den{p}\circ\den{q}$ for all $p,q$. All results
below are \emph{proved} (the $\Set\times\VecC_{\mathrm{fd}}$ result on an
explicitly stated locus).

Everything is controlled by one necessary condition, obtained by evaluating any
candidate composite at the terminal object $\mathbf{1}=[0,X]$.

\begin{lemma}[Shape Lemma; \emph{proved}]\label{lem:cob-shape}
If $\mathcal{C}$ is $\tri$-admissible then for every $P$ and every set $T$ the
internal hom $[P,\,T\cdot\mathbf{1}]$ is a copower of $\mathbf{1}$ (here
$T\cdot\mathbf{1}=\coprod_T\mathbf{1}$). The base must \emph{see} the external
shape set inside itself, as a copower of $\mathbf{1}$.
\end{lemma}

On the set-like extreme admissibility does not merely coexist with the rest of the
structure --- it forces it.

\begin{theorem}[Extensive rigidity and the forced package; \emph{proved}]\label{thm:cob-rigidity}
Let $\mathcal{C}$ be nontrivial, infinitary lextensive and cartesian closed. If
$\mathcal{C}$ is $\tri$-admissible then its unit $\mathbf{1}$ is \emph{connected}
($\mathcal{C}(\mathbf{1},-)$ preserves small coproducts); consequently $\den{-}$
is fully faithful and the reindexing functor $F_q=\Fam(\den{q}^{\op})$ is left
adjoint to $(-)\tri q$ for \emph{every} $q$. In particular $\Set\times\Set$ is
inadmissible, its unit $(1,1)$ being disconnected. Fullness of $\den{-}$ and
existence of the reindexing left adjoint are literally the same comparison lemma
applied twice.
\end{theorem}

The converse fails --- connectedness alone is not enough --- and the exact
condition upgrades the necessary Shape Lemma to a sufficient one.

\begin{theorem}[Characterisation; \emph{proved}]\label{thm:cob-char}
For $\mathcal{C}$ nontrivial, infinitary lextensive and cartesian closed,
$\mathcal{C}$ is $\tri$-admissible iff the Shape Lemma holds, equivalently iff an
external distributive law holds. When $\mathcal{C}$ has a connected-components
functor $\pi_0\dashv(-)\cdot\mathbf{1}$ (as every Grothendieck topos does), this
reads: \emph{$\mathcal{C}$ is $\tri$-admissible iff $\pi_0$ preserves finite
products.}
\end{theorem}

This single criterion unifies the two ways admissibility can fail. The
disconnected unit of $\Set\times\Set$ and the connected-unit obstruction below are
one phenomenon, $\pi_0(X\times Y)\not\cong\pi_0 X\times\pi_0 Y$; it is co-extensive
with the \emph{stable-units} property of the connected-components reflection in
categorical Galois theory \cite{Xarez}, though established here via the
distributive law rather than by descent.

\begin{theorem}[Connected is not sufficient; \emph{proved}]\label{thm:cob-gluing}
The Artin gluing $\mathrm{Gl}((-)^2)$ is a Grothendieck topos with a
\emph{connected} unit that is nonetheless $\tri$-inadmissible: its $\pi_0$ is not
multiplicative, since for the single edge $K$ one has $\pi_0(K\times K)=2\neq1$
(Weichsel's theorem \cite{Weichsel}). This is exactly the gap between the internal
(Gambino--Kock) composition, which every topos supports \cite{GambinoKock}, and
the external-shape composition of $\den{-}$.
\end{theorem}

On the opposite, additive pole the situation inverts. Over $\VecC_{\mathrm{fd}}$
every object is copower-tiny, so $\tri$ collapses onto the Dirichlet tensor,
$p\tri q=p\Day q$, and $(-)\tri q$ has a left adjoint only when $q$ has a single
shape; the unit is always disconnected. Between the poles the admissible region is
thin but inhabited: $\Set\times\VecC_{\mathrm{fd}}$ is admissible (on its locus)
with a disconnected unit, combining a rigid distributive factor with a flexible
copower-tiny one --- so the extensive-pole characterisation does not extend to a
global one (\emph{proved}). The organising fact behind the whole map is
determinacy.

\begin{theorem}[Determinacy; \emph{proved}]\label{thm:cob-determinacy}
If $\mathbf{1}$ is connected then $\den{-}$ is injective on objects up to
isomorphism, so $p\tri q$ is \emph{determined} by $\mathcal{C}$ and admissibility
is a property of $\mathcal{C}$. If $\mathbf{1}$ is disconnected this can fail: over
$\VecC_{\mathrm{fd}}$ the objects $(\{*\},k^2)$ and $(\{1,2\},k)$ present the same
functor $X\mapsto X\oplus X$ but are not isomorphic, so $\tri$ becomes a
\emph{choice} and the base-change question stops being about $\mathcal{C}$ at all.
\end{theorem}

%========================================================================
\section{The linear wall: $\VecC$ is neither full nor faithful}
\label{sec:cob-vec}

The sharpest instrument for separating the four notions of ``container in a
category'' --- external families $\Fam(\mathcal{C}^{\op})$, indexed
dependent-polynomial functors over an LCC base \cite{GambinoKock}, fibrational
containers \cite{vonGlehn}, and Walker's locally subcartesian closed weakening
\cite{Walker} --- is a base that fails \emph{both} extensivity and local cartesian
closure. The category $\VecC$ of vector spaces is exactly such a base: the
comparison $\coprod^{\Set}\to\oplus$ is not an isomorphism, and there is no
internal dependent product $\Pi$. Off the ``good'' zone a base fails for exactly
one of two independent reasons, loss of extensivity (shapes no longer recoverable,
$\den{-}$ no longer full) or loss of internal $\Pi$ ($\tri$ and closedness
collapse); toposes lose neither, quantales and affine bases lose only the second
(reached by Walker's affine repair), and $\VecC$ and $R\text{-}\mathrm{Mod}$ lose
both, so only the external approach reaches them. Three theorems make this
quantitative; all are \emph{proved}.

\begin{theorem}[Fullness $=$ unit-connectedness; \emph{proved}]\label{thm:cob-T1}
For any closed symmetric monoidal $\mathcal{C}$ with small coproducts, $\den{-}$ is
fully faithful iff the monoidal unit $I$ is \emph{connected}
($\mathcal{C}(I,-)$ preserves small coproducts). This is \emph{not} extensivity of
$\mathcal{C}$: $\Set\times\Set$ is (l)extensive yet its unit $(1,1)$ is
disconnected and $\den{-}$ is not full (the two ``crossed'' transformations are
unrealised), while $\Set$ is full not because it is extensive but because $1$ is
connected. It recovers Abbott--Altenkirch--Ghani \cite{AAG} by naming that hidden
hypothesis. The reconstruction of a container from an abstract familially
representable functor is Diers' \cite{Diers}, and uses extensivity of the
\emph{codomain}, not the base --- which is why the folklore said ``extensivity''.
\end{theorem}

\begin{corollary}[The linear collapse; \emph{proved}]\label{cor:cob-veccollapse}
Over $\mathcal{C}=\VecC$ the unit is $k$ and $\VecC(k,-)$ is the underlying-set
functor, which does not preserve $\oplus$: the comparison $\coprod^{\Set}_t
X_t\to\bigoplus_t X_t$ is \emph{neither surjective} (it misses cross-summand sums)
\emph{nor injective} (it collapses the zero vector, which sits in every summand).
Hence over $\VecC$ the extension $\den{-}$ is \emph{neither full nor faithful} ---
the failure of faithfulness being precisely that $\oplus$ collapses tagged zeros.
This is the honest strength of the linear denotation: the on-the-nose
container/functor correspondence that holds over $\Set$ breaks in \emph{both}
directions over an additive base.
\end{corollary}

The other two theorems measure the remaining obstructions, on independent seams.
The Dirichlet tensor $\Day$ is closed iff a familial-representability functor
$\Phi$ is representable; over the linear base this holds \emph{only} on the corner
$\Fam_{\mathrm{fin}}(\VecC_{\mathrm{fd}}^{\op})$ of finite shape-sets and
finite-dimensional positions (the load-bearing conjunct is
\emph{dualizable-and-summable}), failing over both $\Fam(\VecC_{\mathrm{fd}}^{\op})$
and $\Fam(\VecC^{\op})$ by dual mechanisms (\emph{proved}). The substitution
product $\tri$ is left-closed exactly where it collapses onto $\Day$ --- over a
base whose positions are dualizable (tiny), $\tri=\Day$ and the left hom is the
Dirichlet hom --- so it is left-closed on
$\Fam_{\mathrm{fin}}(\VecC_{\mathrm{fd}}^{\op})$ and obstructed over every
extensive base (\emph{proved}). Extensivity is thus \emph{opposed} to
$\tri$-left-closedness: the property whose failure costs fullness and Dirichlet
closedness is exactly what would buy substitution closedness. These are the seams
Weber's distributivity-$\delta$ and familial-functor theories
\cite{WeberPoly,WeberFamilial} place on different rungs; the composition-monoid
refinement of the indexed corner is \cite{DPUV}, and the general-base tensor and
its comonoids-as-enriched-categories are Dorta--Jarvis--Niu \cite{DJN}.

%========================================================================
\section{The $\{\mathrm{Id}\}$ collapse}
\label{sec:cob-collapse}

The linear wall shows a \emph{specific} base failing. The final result is
stronger and structural: over \emph{no} nontrivial effect does the natural
self-enriched semantics enjoy a uniform full-and-faithfulness theorem. The
apparent counterexample --- the reader monad, whose Kleisli category is cartesian
closed --- is the pivot that makes the point.

There are three extensions of an $M$-container into endofunctors, with three
notions of naturality. Fix a monad $M$ and write $\mathcal{C}=\mathrm{Kl}(M)$.

\begin{itemize}
\item[\textbf{(A)}] the \emph{ordinary} extension
  $\den{-}_M\colon\Fam(\mathrm{Kl}(M)^{\op})\to[\Set,\Set]$, whose morphisms are
  pure natural transformations: faithful always, full iff $M$ codense and each
  $\Set(A,M{-})$ connected (Theorem~\ref{thm:cob-codensity});
\item[\textbf{(B)}] the \emph{Kleisli-presheaf} extension
  $\den{-}^{\mathrm{Kl}}\colon\Fam(\mathrm{Kl}(M)^{\op})\to[\mathrm{Kl}(M),\Set]$,
  a coproduct of representables: by ordinary co-Yoneda it is \emph{always} full
  and faithful, at the price that its morphisms are Kleisli-natural;
\item[\textbf{(C)}] the \emph{self-enriched} extension, defined when
  $\mathcal{C}$ is monoidal closed, $\den{S,P}=\sum_s[P_s,-]\colon\mathcal{C}\to
  \mathcal{C}$, which realises the enriched-Yoneda desideratum
  $\Nat(\sum_s[P_s,-],F)\cong\prod_s F(P_s)$ \cite{Kelly}.
\end{itemize}
The three are one factorisation, $\den{-}_M=R\circ\den{-}^{\mathrm{Kl}}$ with
$R=(-)\circ F$ restriction along the free functor $F\colon\Set\to\mathrm{Kl}(M)$;
all loss of ordinary fullness lives in $R$, and its size is exactly the codensity
monad $\Ran_M M$.

Two independent invariants govern (C):
\[
  \mathcal{CL}=\{M:\mathrm{Kl}(M)\text{ monoidal closed}\},\qquad
  \mathcal{UN}=\{M:\text{the monoidal unit of }\mathrm{Kl}(M)\text{ is connected}\}.
\]
Closedness ($\mathcal{CL}$, available by Kock's criterion \cite{KockMonoidal,KockClosed})
makes (C) \emph{exist}; connectedness ($\mathcal{UN}$, i.e.\ $M$ preserves small
coproducts, since $\mathrm{Kl}(M)(1,X)=MX$) makes it \emph{full}
(Theorem~\ref{thm:cob-T1}). They are logically independent --- the cartesian-closed
reader sits in $\mathcal{CL}\setminus\mathcal{UN}$, the writer in $\mathcal{UN}\setminus\mathcal{CL}$ --- and over a
Kleisli base they meet in exactly one monad.

\begin{lemma}[Closed $+$ initial $\Rightarrow$ terminal; \emph{lean-verified}]\label{lem:cob-closedinitial}
In any symmetric monoidal closed category with an initial object $\emptyset$, the
object $[\emptyset,X]$ is terminal for any $X$. (Each $(-)\otimes A$ is a left
adjoint, so $C\otimes\emptyset\cong\emptyset$, whence
$\mathcal{K}(C,[\emptyset,X])\cong\mathcal{K}(\emptyset,X)\cong1$.) The
\emph{existence} of the terminal object uses \emph{only} monoidal closedness and
an initial object --- neither cartesian closedness nor connectedness. The Step-2
\emph{collapse} $E\cong1$ below is another matter: there connectedness
($\mathcal{UN}$, via the writer form) does visible, unavoidable work, namely in
excluding a zero object $\top\cong\emptyset$. Indeed $\mathsf{Maybe}$
($\top\cong\emptyset$, $M\emptyset\cong1$), powerset ($\top\cong\emptyset$), and
the reader ($\top\cong1$) all lie in $\mathcal{CL}$ yet are eliminated only by the
writer form, i.e.\ by membership in $\mathcal{UN}$.
\end{lemma}

\begin{theorem}[$\mathcal{CL}\cap\mathcal{UN}=\{\Id\}$; \emph{proved}, core \emph{lean-verified}]\label{thm:cob-collapse}
Let $M$ be a commutative strong monad on $\Set$. The self-enriched $M$-container
extension $\sum_s[P_s,-]\colon\mathrm{Kl}(M)\to\mathrm{Kl}(M)$ exists and is full
and faithful if and only if $M\cong\Id$. Equivalently $\mathcal{CL}\cap\mathcal{UN}=\{\Id\}$: no nontrivial
effect --- not even an idempotent (localization) monad --- carries a uniform
enriched full-and-faithful semantics.
\end{theorem}

\begin{proof}[Proof sketch]
Fullness of (C) $\iff M\in\mathcal{CL}\cap\mathcal{UN}$. \emph{Step 1:} $\mathcal{UN}$ forces the writer
form. If $M$ preserves small coproducts then the copower-comparison
$T\times M1\to MT$ is a natural bijection, so $M\cong(-)\times E$ with $E=M1$, and
the monad laws are the laws of a (commutative) monoid on $E$. \emph{Step 2:} $\mathcal{CL}$
forces $E\cong1$. The base $\mathrm{Kl}(M)$ is monoidal closed with initial object
$\emptyset$ (since $\mathrm{Kl}(M)(\emptyset,B)=1$), so
Lemma~\ref{lem:cob-closedinitial} supplies a terminal $\top$; then
$\Set(C,M\top)=\mathrm{Kl}(M)(C,\top)\cong1$ forces $M\top\cong1$, and the writer
form gives $|\top|\cdot|E|=1$, so $E\cong1$ and $M\cong\Id$. \emph{Step 3:} $\Id$
qualifies ($\Set$ cartesian closed, $\Set(1,-)=\Id$ preserves coproducts).
\end{proof}

The core of this collapse --- the closed-plus-initial-forces-terminal engine of
Lemma~\ref{lem:cob-closedinitial}, together with the resulting
identity-only classification --- is machine-checked in Lean~4
(\emph{lean-verified}, \texttt{EnrichedFFCollapse}).

\begin{corollary}[No idempotent escape; \emph{proved}]\label{cor:cob-noidem}
No nontrivial idempotent monad lies in $\mathcal{CL}\cap\mathcal{UN}$: membership in $\mathcal{UN}$ forces the
writer form, and a writer monad is idempotent only for $E\cong1$; the terminal
monad $LX=1$ already fails connectedness, $L(1+1)=1\neq2=L1+L1$.
\end{corollary}

\begin{remark}[The reader, reconciled]\label{rem:cob-reader}
The collapse dissolves the apparent paradox. Cartesian closedness of
$\mathrm{Kl}(\mathsf{Reader}_S)$ buys the \emph{availability} of (C), not its
fullness: the unit $1$ has global-points functor $(-)^S$, which fails to preserve
coproducts for $|S|\ge2$, so (C) is not full. The reader's failure is
enrichment-invariant --- non-codensity in (A), a disconnected unit in (C), one
obstruction wearing two masks --- and (B) is full only for Kleisli-natural, not
pure, morphisms. Enrichment relocates the obstruction; it does not remove it.
\end{remark}

\paragraph{The boundary, and the handoff.} The chapter has traced how far the
correctness equation generalises. It survives change of base wholesale --- as
$m$-containers with two computable invariants (\S\ref{sec:cob-mcont}), as the
census of transports (\S\ref{sec:cob-changebase}), as the four monoidal
structures with their closures (\S\ref{sec:cob-four}), and as substitution over
exactly the $\pi_0$-multiplicative bases (\S\ref{sec:cob-admiss}) --- but it does
\emph{not} survive as a \emph{uniform self-enriched full-and-faithful} semantics
over any nontrivial effect: that is the $\{\Id\}$ wall of
Theorem~\ref{thm:cob-collapse}, a statement about general effect monads. A
different wall stands over a \emph{single} additive base: the linear denotation
over $\VecC$ (\S\ref{sec:cob-vec}) is already neither full nor faithful, and the
Vec-admissibility question of whether $F=\Hom(E,\bigoplus_{\N}-)$ is free pushes
the boundary into set theory. That set-theoretic frontier --- where the verdict
becomes independent of ZFC --- is the subject of the next chapter,
Chapter~\ref{ch:frontier}.

% Operators introduced inline (not among the shared macros): $\mathrm{Kl}(M)$
% (Kleisli category), $\mathrm{Gl}$ (Artin gluing), $\pi_0$ (connected components),
% $[-,-]$ (internal hom of the base), $\Phi=\Ran_M M$ (codensity monad),
% $\mathcal{CL},\mathcal{UN}$ (closedness/connectedness classes), $T\cdot\mathbf{1}$ (copower).

% BIBITEMS-USED
% \bibitem{AAG} M.~Abbott, T.~Altenkirch, N.~Ghani, \emph{Containers: constructing strictly positive types}, Theoret.\ Comput.\ Sci.\ \textbf{342} (2005), 3--27.
% \bibitem{AU16} D.~Ahman and T.~Uustalu, \emph{Directed containers as categories}, Electron.\ Proc.\ Theor.\ Comput.\ Sci.\ \textbf{207} (2016), 89--98. arXiv:1604.01187.
% \bibitem{Diers} Y.~Diers, \emph{Cat\'egories localisables}, Th\`ese de doctorat d'\'etat, Universit\'e Paris VI, 1977.
% \bibitem{DJN} A.~Dorta, C.~B.~Jarvis, and N.~Niu, \emph{Monoidal structures on generalized polynomial categories}, arXiv:2305.05655.
% \bibitem{DPUV} G.~De Pascalis, T.~Uustalu, N.~Veltr\`i, \emph{Monoid structures on indexed containers}, arXiv:2509.25879.
% \bibitem{GambinoKock} N.~Gambino and J.~Kock, \emph{Polynomial functors and polynomial monads}, Math.\ Proc.\ Cambridge Philos.\ Soc.\ \textbf{154} (2013), 153--192. arXiv:0906.4931.
% \bibitem{HedgesPoly} J.~Hedges, \emph{A first look at programming in Poly}, blog post, \texttt{julesh.com}, 2026.
% \bibitem{Kelly} G.~M.~Kelly, \emph{Basic Concepts of Enriched Category Theory}, London Math.\ Soc.\ Lecture Note Ser.\ \textbf{64}, Cambridge Univ.\ Press, 1982.
% \bibitem{KockMonoidal} A.~Kock, \emph{Monads on symmetric monoidal closed categories}, Arch.\ Math.\ (Basel) \textbf{21} (1970), 1--10.
% \bibitem{KockClosed} A.~Kock, \emph{Closed categories generated by commutative monads}, J.\ Austral.\ Math.\ Soc.\ \textbf{12} (1971), 405--424.
% \bibitem{NiuSpivak} N.~Niu and D.~I.~Spivak, \emph{Polynomial Functors: A Mathematical Theory of Interaction}, book draft, 2023. arXiv:2312.00990.
% \bibitem{ShapiroSpivak} B.~Shapiro and D.~I.~Spivak, \emph{Structures on categories of polynomials}, arXiv:2305.00167.
% \bibitem{SpivakSrinivasan} D.~I.~Spivak and P.~Srinivasan, \emph{$(\mathbf{Poly},\otimes,\triangleleft)$ is a normal duoidal LDC}, arXiv:2407.01849, 2024.
% \bibitem{vonGlehn} T.~von Glehn, \emph{Polynomials and models of type theory}, PhD thesis, University of Cambridge, 2015.
% \bibitem{Walker} C.~Walker, \emph{Locally subcartesian closed categories}, arXiv:2607.10242.
% \bibitem{WeberPoly} M.~Weber, \emph{Polynomials in categories with pullbacks}, Theory Appl.\ Categ.\ \textbf{30} (2015), 533--598. arXiv:1106.1983.
% \bibitem{WeberFamilial} M.~Weber, \emph{Familial 2-functors and parametric right adjoints}, Theory Appl.\ Categ.\ \textbf{18} (2007), 665--732.
% \bibitem{Weichsel} P.~M.~Weichsel, \emph{The Kronecker product of graphs}, Proc.\ Amer.\ Math.\ Soc.\ \textbf{13} (1962), 47--52.
% \bibitem{Xarez} J.~J.~Xarez, \emph{Admissibility, stable units and connected components}, arXiv:1112.4277, 2011.
